\documentclass[11pt,a4paper]{article}
\usepackage[utf8]{inputenc}
\usepackage[T1]{fontenc}
\usepackage[french]{babel}
\usepackage{amsmath,amssymb,amsfonts,amsthm,mathrsfs,mathtools}
\usepackage[top=2cm,bottom=2cm,left=2cm,right=2cm]{geometry}
\usepackage[dvipsnames]{xcolor}
\usepackage{tcolorbox}
\tcbuselibrary{skins,breakable}
\usepackage{tikz}
\usetikzlibrary{calc,shapes,arrows.meta}
\usepackage{pgfplots}
\pgfplotsset{compat=1.18}
\usepackage{enumitem}
\usepackage{fancyhdr}
\usepackage{lastpage}
\usepackage{hyperref}

% ----------------------------------------------------
% Couleurs officielles du projet (Charte LLG / MPSI)
% ----------------------------------------------------
\definecolor{llgblue}{RGB}{20, 50, 115}
\definecolor{llglight}{RGB}{246, 249, 255}

\definecolor{emeraldframe}{RGB}{18, 110, 75}
\definecolor{emeraldback}{RGB}{242, 250, 246}
\definecolor{emeraldtitle}{RGB}{12, 80, 52}

\definecolor{prepaframe}{RGB}{30, 70, 135}
\definecolor{prepaback}{RGB}{246, 249, 255}
\definecolor{prepatitle}{RGB}{20, 50, 105}

\definecolor{alertframe}{RGB}{185, 75, 20}
\definecolor{alertback}{RGB}{255, 249, 245}
\definecolor{alerttitle}{RGB}{145, 55, 12}

% Configuration des hyperliens
\hypersetup{
    colorlinks=true,
    linkcolor=llgblue,
    urlcolor=llgblue,
    citecolor=llgblue
}

% En-têtes et pieds de page
\pagestyle{fancy}
\fancyhf{}
\renewcommand{\headrulewidth}{0.5pt}
\renewcommand{\footrulewidth}{0.4pt}
\fancyhead[L]{\small\textbf{Lycée Louis-le-Grand} -- 1\textsuperscript{ère} Spécialité Mathématiques (1E3)}
\fancyhead[C]{\small\textbf{DST : Fonction Exponentielle}}
\fancyhead[R]{\small 8 Mars 2022}
\fancyfoot[L]{\footnotesize LLG 2021/2022 -- Devoir Sur Table (1E3)}
\fancyfoot[C]{\footnotesize Énoncé \& Corrigé Intégral Rigoureux}
\fancyfoot[R]{\footnotesize Page \thepage\ / \pageref{LastPage}}

% ----------------------------------------------------
% Boîtes tcolorbox personnalisées
% ----------------------------------------------------
\newtcolorbox{correctionbox}[1][]{
  colback=emeraldback,
  colframe=emeraldframe,
  fonttitle=\bfseries\color{white},
  coltitle=white,
  colbacktitle=emeraldtitle,
  title={Corrigé officiel détaillé et barème indicatif #1},
  arc=2mm,
  boxrule=0.8pt,
  titlerule=0.5pt,
  breakable,
  before skip=0.4cm,
  after skip=0.5cm
}

\newtcolorbox{remarquebox}[1][]{
  colback=prepaback,
  colframe=prepaframe,
  fonttitle=\bfseries\small\color{white},
  coltitle=white,
  colbacktitle=prepatitle,
  title={Prolongement \& Regard Prépa / MPSI #1},
  arc=1.5mm,
  boxrule=0.6pt,
  titlerule=0.4pt,
  breakable,
  before skip=0.25cm,
  after skip=0.35cm
}

\newtcolorbox{alertebox}[1][]{
  colback=alertback,
  colframe=alertframe,
  fonttitle=\bfseries\small\color{white},
  coltitle=white,
  colbacktitle=alerttitle,
  title={Point de vigilance méthodologique \& Pièges fréquents #1},
  arc=1.5mm,
  boxrule=0.6pt,
  titlerule=0.4pt,
  breakable,
  before skip=0.25cm,
  after skip=0.35cm
}

% Commandes mathématiques usuelles
\newcommand{\N}{\mathbb{N}}
\newcommand{\Z}{\mathbb{Z}}
\newcommand{\Q}{\mathbb{Q}}
\newcommand{\R}{\mathbb{R}}
\newcommand{\C}{\mathbb{C}}
\newcommand{\dd}{\mathrm{d}}
\newcommand{\e}{\mathrm{e}}

\begin{document}

% ====================================================
% EN-TÊTE OFFICIEL DU DEVOIR
% ====================================================
\begin{center}
    {\LARGE \bfseries Lycée Louis-le-Grand -- Paris}\\[0.2cm]
    {\large \textbf{Classe de Première Spécialité Mathématiques -- Classe 1E3}}\\[0.15cm]
    {\normalsize Année scolaire 2021--2022}\\[0.35cm]
    \hrule height 1.2pt
    \vspace{0.25cm}
    {\LARGE \bfseries DEVOIR SUR TABLE : FONCTION EXPONENTIELLE}\\[0.15cm]
    {\large \textbf{Algèbre rationnelle, Dérivée seconde, Point d'inflexion, Distance minimale \& Convexité}}\\[0.2cm]
    \hrule height 0.6pt
    \vspace{0.25cm}
    {\small \textbf{Date :} Mardi 8 Mars 2022 \quad $\vert$ \quad \textbf{Durée :} 2 heures \quad $\vert$ \quad \textbf{Calculatrice autorisée (approximations)}}\\[0.1cm]
    {\small \textbf{Barème total :} 20 points \quad (Ex 1 : 2 pts \quad $\vert$ \quad Ex 2 : 7 pts \quad $\vert$ \quad Ex 3 : 9 pts \quad $\vert$ \quad Ex 4 : 2 pts)}
\end{center}

\vspace{0.3cm}

\begin{tcolorbox}[colback=llglight,colframe=llgblue,arc=1.5mm,boxrule=0.7pt]
\textbf{Consignes officielles d'évaluation :}
\begin{itemize}[leftmargin=*,itemsep=2pt]
    \item La clarté, la précision des justifications et la rigueur de la rédaction constituent des critères fondamentaux de notation.
    \item Les théorèmes utilisés doivent être explicitement mentionnés (théorèmes généraux de dérivation, dérivation des composées, corollaire du théorème des valeurs intermédiaires pour les fonctions strictement monotones).
    \item L'usage de la calculatrice est strictement restreint aux calculs d'approximations numériques de l'exercice 3.
    \item Tous les résultats finaux doivent être encadrés avec soin.
\end{itemize}
\end{tcolorbox}

\vspace{0.4cm}

% ####################################################
% PARTIE I : ÉNONCÉS INTÉGRAUX (SUJET DU DST)
% ####################################################
\begin{center}
    \rule{0.35\textwidth}{0.8pt}\\[0.15cm]
    {\Large \bfseries PREMIÈRE PARTIE : ÉNONCÉ DU SUJET}\\[0.1cm]
    \rule{0.35\textwidth}{0.8pt}
\end{center}

\vspace{0.2cm}

% ----------------------------------------------------
% EXERCICE 1 (ÉNONCÉ)
% ----------------------------------------------------
\subsection*{Exercice 1 \hfill \normalfont\normalsize (2 points)}
\textbf{Simplification algébrique rationnelle en exponentielles}

\medskip
Soit $A$ le nombre réel défini pour tout réel $x$ par l'expression :
\[
A = 1 - \frac{\e^x}{\e^x + 1} - \frac{\e^x}{(\e^x + 1)^2}.
\]

Exprimer $A$ sous la forme d'un quotient dont le numérateur est égal à $1$.

\vspace{0.5cm}

% ----------------------------------------------------
% EXERCICE 2 (ÉNONCÉ)
% ----------------------------------------------------
\subsection*{Exercice 2 \hfill \normalfont\normalsize (7 points)}
\textbf{Tangente, dérivée seconde et point d'inflexion}

\medskip
On considère la fonction $f$ définie sur $\R$ par :
\[
f(x) = \frac{x+2}{\e^x}.
\]
On note $\mathcal{C}_f$ sa courbe représentative dans un repère orthogonal du plan.

\begin{enumerate}[label=\textbf{\arabic*.}]
    \item Déterminer l'équation réduite de la tangente $T$ à la courbe $\mathcal{C}_f$ au point d'abscisse $0$. \hfill \textbf{(2 points)}
    \item On considère la fonction auxiliaire $g$ définie sur $\R$ par :
    \[
    g(x) = x + 2 + (x-2)\e^x.
    \]
    \begin{enumerate}[label=\textbf{(\alph*)}]
        \item Donner l'expression de $g''(x)$ pour tout réel $x$, où $g''$ désigne la fonction dérivée seconde de la fonction $g$. \hfill \textbf{(1{,}5 point)}
        \item En déduire le tableau de variation de la fonction $g$ sur $\R$. On prendra soin de calculer les valeurs exactes de $g(0)$ et de $g'(0)$. \hfill \textbf{(2 points)}
    \end{enumerate}
    \item Étudier la position relative de la tangente $T$ par rapport à la courbe $\mathcal{C}_f$. Interpréter géométriquement la nature du point de contact d'abscisse $0$. \hfill \textbf{(1{,}5 point)}
\end{enumerate}

\vspace{0.5cm}

% ----------------------------------------------------
% EXERCICE 3 (ÉNONCÉ)
% ----------------------------------------------------
\subsection*{Exercice 3 \hfill \normalfont\normalsize (9 points)}
\textbf{Distance minimale à la courbe de l'exponentielle et orthogonalité normale}

\medskip
\textbf{Partie A : Étude d'une fonction auxiliaire}

On considère la fonction $g$ définie sur $\R$ par :
\[
g(x) = x + \e^{2x}.
\]
\begin{enumerate}[label=\textbf{\arabic*.}]
    \item Déterminer les limites de la fonction $g$ en $-\infty$ et en $+\infty$. \hfill \textbf{(1 point)}
    \item Déterminer l'expression de la fonction dérivée $g'(x)$ pour tout réel $x$. \hfill \textbf{(0{,}5 point)}
    \item Dresser le tableau de variation complet de la fonction $g$ sur $\R$ en y faisant figurer les limites obtenues. \hfill \textbf{(1 point)}
    \item On admet que l'équation $g(x) = 0$ admet une unique solution réelle sur $\R$, notée $\alpha$. En déduire le tableau de signe de $g(x)$ sur $\R$ en faisant figurer $\alpha$. \hfill \textbf{(1 point)}
    \item À l'aide de la calculatrice, donner une valeur approchée de $\alpha$ arrondie à $10^{-1}$ près. \hfill \textbf{(0{,}5 point)}
\end{enumerate}

\medskip
\textbf{Partie B : Distance minimale à l'origine et propriété géométrique}

On considère la fonction $f$ définie sur $\R$ par :
\[
f(x) = \sqrt{x^2 + \e^{2x}}.
\]
\begin{enumerate}[label=\textbf{\arabic*.}]
    \item Justifier que la fonction $f$ est dérivable sur $\R$ et montrer que pour tout réel $x$ :
    \[
    f'(x) = \frac{g(x)}{\sqrt{x^2 + \e^{2x}}}.
    \] \hfill \textbf{(1{,}5 point)}
    \item On se place dans le plan muni d'un repère orthonormé d'origine $O$. On note $\mathcal{C}_{\exp}$ la courbe représentative de la fonction exponentielle standard ($y = \e^x$).
    \begin{enumerate}[label=\textbf{(\alph*)}]
        \item Déterminer, en fonction du réel $\alpha$ défini dans la Partie A, les coordonnées du point $A$ de la courbe $\mathcal{C}_{\exp}$ situé à distance minimale de l'origine $O$. \hfill \textbf{(2 points)}
        \item Démontrer qu'au point $A$, la tangente $T_A$ à la courbe de l'exponentielle est rigoureusement perpendiculaire à la droite $(OA)$. \hfill \textbf{(1{,}5 point)}
    \end{enumerate}
\end{enumerate}

\vspace{0.5cm}

% ----------------------------------------------------
% EXERCICE 4 (ÉNONCÉ)
% ----------------------------------------------------
\subsection*{Exercice 4 \hfill \normalfont\normalsize (2 points)}
\textbf{Inégalité de moyennes et convexité de la fonction exponentielle}

\medskip
Soient $x$ et $y$ deux nombres réels quelconques.

Comparer les deux quantités réelles :
\[
\e^{\frac{x+y}{2}} \quad \text{et} \quad \frac{\e^x + \e^y}{2}.
\]
Préciser la condition nécessaire et suffisante pour que l'égalité soit réalisée.

\newpage

% ####################################################
% PARTIE II : CORRIGÉS DÉTAILLÉS ET APPROFONDISSEMENTS
% ####################################################
\begin{center}
    \rule{0.35\textwidth}{0.8pt}\\[0.15cm]
    {\Large \bfseries SECONDE PARTIE : CORRIGÉ OFFICIEL \& REGARD PRÉPA}\\[0.1cm]
    \rule{0.35\textwidth}{0.8pt}
\end{center}

\vspace{0.3cm}

% ====================================================
% CORRIGÉ EXERCICE 1
% ====================================================
\subsection*{Exercice 1 : Simplification algébrique rationnelle \hfill \normalfont\normalsize (2 points)}

\begin{correctionbox}
\textbf{1. Regroupement et réduction au même dénominateur des deux premiers termes :}

Pour tout nombre réel $x$, $\e^x > 0$, donc $\e^x + 1 > 1 \neq 0$. L'expression $A$ est donc parfaitement définie sur $\R$.
Effectuons d'abord la soustraction des deux premiers termes en réduisant au dénominateur commun $(\e^x + 1)$ :
\[
1 - \frac{\e^x}{\e^x + 1} = \frac{(\e^x + 1) - \e^x}{\e^x + 1} = \frac{1}{\e^x + 1}.
\] \hfill \textbf{[1 point]}

\medskip
\textbf{2. Soustraction du troisième terme et factorisation finale :}

En substituant ce résultat dans l'expression complète de $A$, on obtient :
\[
A = \frac{1}{\e^x + 1} - \frac{\e^x}{(\e^x + 1)^2}.
\]
Le dénominateur commun est désormais $(\e^x + 1)^2$. En multipliant le numérateur et le dénominateur de la première fraction par $(\e^x + 1)$ :
\[
A = \frac{\e^x + 1}{(\e^x + 1)^2} - \frac{\e^x}{(\e^x + 1)^2} = \frac{(\e^x + 1) - \e^x}{(\e^x + 1)^2} = \frac{1}{(\e^x + 1)^2}.
\]
L'expression de $A$ sous forme d'un quotient dont le numérateur vaut $1$ est donc :
\[
\mathbf{A = \frac{1}{(\e^x + 1)^2}}.
\] \hfill \textbf{[1 point]}
\end{correctionbox}

\begin{alertebox}[Stratégie de calcul : Pas à pas vs Dénominateur global]
Une erreur fréquente d'élèves consiste à réduire d'emblée l'ensemble des trois termes au dénominateur commun $(\e^x + 1)^2$ :
\[
A = \frac{(\e^x + 1)^2 - \e^x(\e^x + 1) - \e^x}{(\e^x + 1)^2} = \frac{(\e^{2x} + 2\e^x + 1) - (\e^{2x} + \e^x) - \e^x}{(\e^x + 1)^2}.
\]
Bien que tout à fait exacte, cette démarche présente un risque accru d'erreur de signe lors de la distribution du signe moins : $-(\e^{2x} + \e^x) = -\e^{2x} - \e^x$. La décomposition en cascade adoptée dans le corrigé est à la fois plus rapide, plus élégante et immunisée contre ces pièges élémentaires.
\end{alertebox}

\begin{remarquebox}[Lien avec la fonction sigmoïde et dérivation formelle]
D'un point de vue supérieur, l'expression étudiée est intimement liée à la \textbf{fonction logistique standard} (ou sigmoïde), omniprésente en probabilités, en apprentissage automatique (\textit{machine learning}) et en dynamique des populations :
\[
\sigma(x) = \frac{1}{1 + \e^{-x}} = \frac{\e^x}{\e^x + 1}.
\]
On observe les propriétés fondamentales suivantes :
\begin{enumerate}
    \item $1 - \sigma(x) = 1 - \frac{\e^x}{\e^x + 1} = \frac{1}{\e^x + 1}$.
    \item La dérivée de la sigmoïde vérifie la relation différentielle remarquable :
    \[
    \sigma'(x) = \frac{\e^x(\e^x + 1) - \e^x(\e^x)}{(\e^x + 1)^2} = \frac{\e^x}{(\e^x + 1)^2} = \sigma(x)\bigl(1 - \sigma(x)\bigr).
    \]
\end{enumerate}
L'expression $A$ de l'exercice se réécrit donc sous forme condensée :
\[
A = 1 - \sigma(x) - \sigma'(x) = \bigl(1 - \sigma(x)\bigr) - \sigma(x)\bigl(1 - \sigma(x)\bigr) = \bigl(1 - \sigma(x)\bigr)^2 = \left(\frac{1}{\e^x + 1}\right)^2 = \frac{1}{(\e^x + 1)^2}.
\]
L'identité algébrique découle ainsi directement de l'identité polynomiale élémentaire $1 - u - u(1 - u) = (1 - u)^2$ posée avec $u = \sigma(x)$.
\end{remarquebox}

\vspace{0.4cm}

% ====================================================
% CORRIGÉ EXERCICE 2
% ====================================================
\subsection*{Exercice 2 : Tangente, dérivée seconde et point d'inflexion \hfill \normalfont\normalsize (7 points)}

\begin{correctionbox}
\textbf{1. Équation réduite de la tangente $T$ au point d'abscisse $0$ :}

La fonction $f : x \mapsto \frac{x+2}{\e^x}$ est définie sur $\R$.
Les fonctions $u : x \mapsto x + 2$ et $v : x \mapsto \e^x$ sont dérivables sur $\R$, avec pour tout $x \in \R$, $v(x) = \e^x \neq 0$.
Par théorème sur le quotient de fonctions dérivables, $f$ est dérivable sur $\R$. Pour tout réel $x$ :
\[
f'(x) = \frac{u'(x)v(x) - u(x)v'(x)}{\bigl(v(x)\bigr)^2} = \frac{1 \cdot \e^x - (x+2)\e^x}{(\e^x)^2} = \frac{\e^x\bigl(1 - (x + 2)\bigr)}{(\e^x)^2} = \frac{-x - 1}{\e^x}.
\] \hfill \textbf{[1 point]}

Calculons l'image de $0$ par $f$ et par sa dérivée $f'$ :
\[
f(0) = \frac{0 + 2}{\e^0} = \frac{2}{1} = 2,
\]
\[
f'(0) = \frac{-0 - 1}{\e^0} = \frac{-1}{1} = -1.
\]
L'équation réduite de la tangente $T$ au point d'abscisse $0$ s'écrit classiquement :
\[
y = f'(0)(x - 0) + f(0) \iff y = -1 \cdot x + 2 \iff \mathbf{y = -x + 2}.
\] \hfill \textbf{[1 point]}

\medskip
\textbf{2. (a) Expression de la dérivée seconde $g''(x)$ :}

La fonction $g$ est définie sur $\R$ par $g(x) = x + 2 + (x-2)\e^x$.
Elle est indéfiniment dérivable sur $\R$ comme somme et produit de fonctions usuelles.
Calculons d'abord sa dérivée première $g'$ :
Pour tout $x \in \R$ :
\[
g'(x) = 1 + 0 + \left[ \frac{\dd}{\dd x}(x - 2) \cdot \e^x + (x - 2) \cdot \frac{\dd}{\dd x}(\e^x) \right] = 1 + 1 \cdot \e^x + (x - 2)\e^x = 1 + (1 + x - 2)\e^x,
\]
\[
g'(x) = 1 + (x - 1)\e^x.
\] \hfill \textbf{[0{,}75 point]}

Dérivons à nouveau pour obtenir la fonction dérivée seconde $g''$ :
\[
g''(x) = 0 + \left[ \frac{\dd}{\dd x}(x - 1) \cdot \e^x + (x - 1) \cdot \frac{\dd}{\dd x}(\e^x) \right] = 1 \cdot \e^x + (x - 1)\e^x = (1 + x - 1)\e^x,
\]
\[
\mathbf{g''(x) = x\e^x}.
\] \hfill \textbf{[0{,}75 point]}

\medskip
\textbf{2. (b) Tableau de variation de la fonction $g$ :}

\begin{itemize}[leftmargin=*]
    \item \textbf{Signe de la dérivée seconde $g''(x)$ et variations de $g'$ :}
    Pour tout réel $x$, la fonction exponentielle vérifie $\e^x > 0$.
    Par conséquent, le signe de $g''(x) = x\e^x$ est strictement identique à celui de $x$ :
    \begin{itemize}
        \item Pour tout $x \in ]-\infty\,;\,0[$, $g''(x) < 0$, donc $g'$ est strictement décroissante sur $]-\infty\,;\,0]$.
        \item Pour tout $x \in ]0\,;\,+\infty[$, $g''(x) > 0$, donc $g'$ est strictement croissante sur $[0\,;\,+\infty[$.
        \item En $x = 0$, $g''(0) = 0$.
    \end{itemize}
    \item \textbf{Calcul de l'extremum de $g'$ en $x = 0$ :}
    La fonction $g'$ admet ainsi un minimum global strict en $x = 0$ sur $\R$. La valeur de ce minimum est :
    \[
    g'(0) = 1 + (0 - 1)\e^0 = 1 - 1 = \mathbf{0}.
    \] \hfill \textbf{[0{,}5 point]}
    \item \textbf{Signe de $g'(x)$ et sens de variation de $g$ :}
    Puisque le minimum absolu de $g'$ sur $\R$ est nul et qu'il n'est atteint qu'en l'unique point isolé $x = 0$, on a :
    \[
    \forall x \in \R \setminus \{0\}, \quad g'(x) > g'(0) = 0 \quad \text{et} \quad g'(0) = 0.
    \]
    La dérivée première $g'$ est donc strictement positive sur $\R^*$.
    On en déduit que la fonction $g$ est \textbf{strictement croissante sur $\R$}. \hfill \textbf{[0{,}75 point]}
    \item \textbf{Calcul de $g(0)$ et déduction du signe de $g$ :}
    Calculons l'image de $0$ par la fonction $g$ :
    \[
    g(0) = 0 + 2 + (0 - 2)\e^0 = 2 - 2 = \mathbf{0}.
    \]
    Puisque $g$ est strictement croissante sur $\R$ et s'annule en $0$, on en déduit immédiatement son tableau de signe :
    \begin{itemize}
        \item Pour tout $x \in ]-\infty\,;\,0[$, $\mathbf{g(x) < 0}$ ;
        \item Pour $x = 0$, $\mathbf{g(0) = 0}$ ;
        \item Pour tout $x \in ]0\,;\,+\infty[$, $\mathbf{g(x) > 0}$.
    \end{itemize} \hfill \textbf{[0{,}75 point]}
\end{itemize}

\begin{center}
\begin{tikzpicture}[scale=0.9]
    \draw[thick] (0,0) rectangle (12,4.5);
    \draw[thick] (0,1.5) -- (12,1.5);
    \draw[thick] (0,3) -- (12,3);
    \draw[thick] (3,0) -- (3,4.5);
    
    \node at (1.5,3.75) {$x$};
    \node at (3.5,3.75) {$-\infty$};
    \node at (7.5,3.75) {$0$};
    \node at (11.5,3.75) {$+\infty$};
    
    \node at (1.5,2.25) {$g'(x)$};
    \node at (5.5,2.5) {};
    \node at (7.5,2.25) {$0$ (min)};
    \node at (5.5,1.9) {$+$};
    \node at (9.5,1.9) {$+$};
    
    \node at (1.5,0.75) {$g(x)$};
    \node at (3.5,0.3) {};
    \node at (7.5,0.75) {$0$};
    \draw[->, thick, >=Latex] (4,0.3) -- (11,1.2);
\end{tikzpicture}
\end{center}

\medskip
\textbf{3. Position relative de $T$ et $\mathcal{C}_f$ et interprétation géométrique :}

Pour étudier la position relative de la courbe $\mathcal{C}_f$ par rapport à sa tangente $T$, étudions le signe de la différence $d(x) = f(x) - (-x + 2)$ pour tout réel $x$ :
\[
d(x) = \frac{x+2}{\e^x} - (-x + 2) = \frac{x + 2 - (-x + 2)\e^x}{\e^x} = \frac{x + 2 + (x - 2)\e^x}{\e^x}.
\]
On reconnaît au numérateur l'expression exacte de la fonction auxiliaire $g(x)$ :
\[
d(x) = \frac{g(x)}{\e^x}.
\] \hfill \textbf{[0{,}5 point]}

Comme la fonction exponentielle est strictement positive sur $\R$ ($\e^x > 0$), le signe de la différence $d(x)$ est strictement identique à celui de $g(x)$ :
\begin{itemize}[leftmargin=*]
    \item Sur l'intervalle $]-\infty\,;\,0[$, $g(x) < 0 \implies d(x) < 0$.
    La courbe $\mathcal{C}_f$ est donc \textbf{strictement située en-dessous de la tangente $T$}.
    \item Sur l'intervalle $]0\,;\,+\infty[$, $g(x) > 0 \implies d(x) > 0$.
    La courbe $\mathcal{C}_f$ est donc \textbf{strictement située au-dessus de la tangente $T$}.
    \item En $x = 0$, $d(0) = 0$. La courbe $\mathcal{C}_f$ et sa tangente $T$ ont un unique point d'intersection commun : le point de coordonnées $I(0\,;\,2)$.
\end{itemize}
\textbf{Interprétation géométrique :} La droite $T$, tangente à la courbe au point d'abscisse $0$, traverse la courbe $\mathcal{C}_f$ en ce point en changeant de position relative. Par définition, le point de contact $\mathbf{I(0\,;\,2)}$ est un \textbf{point d'inflexion} de la courbe représentative $\mathcal{C}_f$. \hfill \textbf{[1 point]}
\end{correctionbox}

\begin{center}
\begin{tikzpicture}[scale=1.2, cap=round, >=Latex]
    % Grille et axes
    \draw[very thin, gray!30] (-2.5,-1) grid (3.5,4.5);
    \draw[->, thick] (-2.5,0) -- (3.8,0) node[right] {\small $x$};
    \draw[->, thick] (0,-1) -- (0,4.8) node[above] {\small $y$};
    \node[below left] at (0,0) {\small $O$};
    \draw[thick] (1,0.08) -- (1,-0.08) node[below] {\small $1$};
    \draw[thick] (0.08,1) -- (-0.08,1) node[left] {\small $1$};

    % Courbe f(x) = (x+2)*exp(-x)
    \draw[domain=-1.5:3.5, samples=150, very thick, llgblue] 
        plot (\x, {(\x + 2)*exp(-\x)}) 
        node[above right] {\small $\mathcal{C}_f : y = \dfrac{x+2}{\e^x}$};

    % Tangente T : y = -x + 2
    \draw[domain=-1.2:3.2, samples=50, thick, alertframe, dashed] 
        plot (\x, {-\x + 2}) 
        node[right] {\small $T : y = -x + 2$};

    % Point d'inflexion I(0, 2)
    \filldraw[ForestGreen] (0,2) circle (2.5pt) node[right=4pt] {\small $I(0\,;\,2)$ \textbf{(Point d'inflexion)}};

    % Zones de position relative
    \node[llgblue, fill=white, inner sep=1pt] at (-1.1, 1.8) {\footnotesize $\mathcal{C}_f$ en-dessous de $T$};
    \node[llgblue, fill=white, inner sep=1pt] at (2.2, 1.2) {\footnotesize $\mathcal{C}_f$ au-dessus de $T$};

    % Maximum local en x = -1 : f(-1) = e approx 2.718
    \filldraw[gray!70] (-1, 2.718) circle (1.8pt);
    \draw[<->, dashed, gray!70] (-1.5, 2.718) -- (-0.5, 2.718) node[above right] {\footnotesize Sommet $(-1\,;\,\e)$};
\end{tikzpicture}
\end{center}

\begin{remarquebox}[Lien direct avec la convexité de $f$ et formule de Taylor]
D'un point de vue supérieur (programme de Terminale / MPSI), la détermination d'un point d'inflexion ne nécessite pas d'introduire artificiellement la fonction $g$. Elle découle directement du calcul de la dérivée seconde $f''$ de la fonction $f$.
Reprenons $f'(x) = (-x - 1)\e^{-x}$. Par dérivation d'un produit :
\[
f''(x) = \frac{\dd}{\dd x}(-x - 1) \cdot \e^{-x} + (-x - 1) \cdot \frac{\dd}{\dd x}(\e^{-x}) = -1 \cdot \e^{-x} + (-x - 1)(-\e^{-x}) = (-\e^{-x}) + (x + 1)\e^{-x} = x\e^{-x} = \frac{x}{\e^x}.
\]
Puisque $\e^x > 0$, la dérivée seconde $f''(x)$ s'annule en $x = 0$ en changeant de signe :
\begin{itemize}
    \item Pour $x < 0$, $f''(x) < 0$ : la fonction $f$ est \textbf{strictement concave} sur $]-\infty\,;\,0]$, sa courbe se situe donc sous toutes ses tangentes.
    \item Pour $x > 0$, $f''(x) > 0$ : la fonction $f$ est \textbf{strictement convexe} sur $[0\,;\,+\infty[$, sa courbe se situe donc au-dessus de toutes ses tangentes.
\end{itemize}
Le changement de convexité en $x = 0$ démontre instantanément que $(0\,;\,2)$ est un point d'inflexion.
Par ailleurs, le développement de Taylor au voisinage de $0$ donne :
\[
f(x) = (x+2)\left(1 - x + \frac{x^2}{2} - \frac{x^3}{6} + o(x^3)\right) = 2 - x + 0\cdot x^2 + \frac{x^3}{3} + o(x^3).
\]
La partie affine $y = 2 - x$ est bien la tangente $T$, et le premier terme d'écart non nul est $\frac{x^3}{3}$, qui est de degré impair (ordre 3), traduisant algébriquement la traversée de la tangente par la courbe.
\end{remarquebox}

\vspace{0.4cm}

% ====================================================
% CORRIGÉ EXERCICE 3
% ====================================================
\subsection*{Exercice 3 : Distance minimale et orthogonalité normale \hfill \normalfont\normalsize (9 points)}

\begin{correctionbox}
\textbf{Partie A : Étude d'une fonction auxiliaire}

\textbf{1. Limites de la fonction $g$ aux bornes de $\R$ :}

La fonction $g$ est définie sur $\R$ par $g(x) = x + \e^{2x}$.
\begin{itemize}[leftmargin=*]
    \item \textbf{En $-\infty$ :} On a $\lim_{x \to -\infty} x = -\infty$.
    Posons $X = 2x$. Lorsque $x \to -\infty$, $X \to -\infty$, donc $\lim_{x \to -\infty} \e^{2x} = \lim_{X \to -\infty} \e^X = 0$.
    Par somme des limites :
    \[
    \lim_{x \to -\infty} g(x) = -\infty.
    \] \hfill \textbf{[0{,}5 point]}
    \item \textbf{En $+\infty$ :} On a $\lim_{x \to +\infty} x = +\infty$ et $\lim_{x \to +\infty} \e^{2x} = +\infty$.
    Par somme de limites de termes divergeant vers $+\infty$, il n'y a aucune indétermination :
    \[
    \lim_{x \to +\infty} g(x) = +\infty.
    \] \hfill \textbf{[0{,}5 point]}
\end{itemize}

\medskip
\textbf{2. Expression de la fonction dérivée $g'(x)$ :}

La fonction $g$ est dérivable sur $\R$ comme somme de la fonction affine $x \mapsto x$ et de la fonction composée $x \mapsto \e^{2x}$ ($x \mapsto \e^{u(x)}$ avec $u(x) = 2x$).
Pour tout réel $x$ :
\[
g'(x) = 1 + u'(x)\e^{u(x)} = 1 + 2\e^{2x}.
\] \hfill \textbf{[0{,}5 point]}

\medskip
\textbf{3. Tableau de variation de $g$ sur $\R$ :}

Pour tout $x \in \R$, on sait que la fonction exponentielle est strictement positive, donc $\e^{2x} > 0$, d'où $2\e^{2x} > 0$.
On en déduit que pour tout réel $x$ :
\[
g'(x) = 1 + 2\e^{2x} > 1 > 0.
\]
La dérivée $g'$ étant strictement positive sur $\R$, la fonction $g$ est \textbf{strictement croissante sur $\R$}. \hfill \textbf{[1 point]}

\begin{center}
\begin{tikzpicture}[scale=0.9]
    \draw[thick] (0,0) rectangle (10,3);
    \draw[thick] (0,1.5) -- (10,1.5);
    \draw[thick] (2.5,0) -- (2.5,3);
    
    \node at (1.25,2.25) {$x$};
    \node at (3,2.25) {$-\infty$};
    \node at (6.25,2.25) {$\alpha$};
    \node at (9.5,2.25) {$+\infty$};
    
    \node at (1.25,0.75) {$g(x)$};
    \node at (3,0.3) {$-\infty$};
    \node at (6.25,0.75) {$0$};
    \node at (9.5,1.2) {$+\infty$};
    \draw[->, thick, >=Latex] (3.4,0.3) -- (9.1,1.2);
\end{tikzpicture}
\end{center}

\medskip
\textbf{4. Déduction du tableau de signe de $g(x)$ :}

La fonction $g$ est continue (car dérivable) et strictement croissante sur $\R$.
Elle réalise une bijection de $\R$ sur $]\lim_{-\infty} g\,;\, \lim_{+\infty} g[ = ]-\infty\,;\,+\infty[ = \R$.
Comme $0 \in \R$, il existe un unique antécédent réel de $0$, noté $\alpha$, tel que $g(\alpha) = 0$.
Par stricte croissance de $g$ sur $\R$ :
\begin{itemize}[leftmargin=*]
    \item Pour tout $x \in ]-\infty\,;\,\alpha[$, $g(x) < g(\alpha)$, soit $\mathbf{g(x) < 0}$ ;
    \item Pour $x = \alpha$, $\mathbf{g(\alpha) = 0}$ ;
    \item Pour tout $x \in ]\alpha\,;\,+\infty[$, $g(x) > g(\alpha)$, soit $\mathbf{g(x) > 0}$.
\end{itemize} \hfill \textbf{[1 point]}

\begin{center}
\begin{tikzpicture}[scale=0.9]
    \draw[thick] (0,0) rectangle (10,1.5);
    \draw[thick] (2.5,0) -- (2.5,1.5);
    
    \node at (1.25,0.75) {$x$};
    \node at (3,0.75) {$-\infty$};
    \node at (6.25,0.75) {$\alpha$};
    \node at (9.5,0.75) {$+\infty$};
    
    \draw[thick] (0,-1) rectangle (10,0);
    \draw[thick] (2.5,-1) -- (2.5,0);
    \node at (1.25,-0.5) {$g(x)$};
    \node at (4.5,-0.5) {$-$};
    \node at (6.25,-0.5) {$0$};
    \node at (8,-0.5) {$+$};
\end{tikzpicture}
\end{center}

\medskip
\textbf{5. Approximation numérique de $\alpha$ :}

À l'aide de la calculatrice, évaluons $g(x)$ au pas de $0{,}1$ :
\[
g(-0{,}5) = -0{,}5 + \e^{2(-0{,}5)} = -0{,}5 + \e^{-1} \approx -0{,}5 + 0{,}3679 = -0{,}1321 < 0,
\]
\[
g(-0{,}4) = -0{,}4 + \e^{2(-0{,}4)} = -0{,}4 + \e^{-0{,}8} \approx -0{,}4 + 0{,}4493 = +0{,}0493 > 0.
\]
Puisque $g(-0{,}5) < 0 < g(-0{,}4)$, le théorème des valeurs intermédiaires garantit que $-0{,}5 < \alpha < -0{,}4$.
En affinant la recherche (dichotomie ou balayage au centième) :
\[
g(-0{,}43) \approx -0{,}43 + \e^{-0{,}86} \approx -0{,}0068 < 0 \quad \text{et} \quad g(-0{,}42) \approx -0{,}42 + \e^{-0{,}84} \approx +0{,}0117 > 0,
\]
ce qui donne $\alpha \approx -0{,}426$. Arrondi à $10^{-1}$ près, on a :
\[
\mathbf{\alpha \approx -0{,}4}.
\] \hfill \textbf{[0{,}5 point]}

\vspace{0.3cm}
\hrule
\vspace{0.3cm}

\textbf{Partie B : Distance minimale à l'origine et orthogonalité}

\textbf{1. Dérivabilité et expression de $f'(x)$ :}

Soit la fonction $u : x \mapsto x^2 + \e^{2x}$.
Pour tout réel $x$, $x^2 \ge 0$ et $\e^{2x} > 0$, donc la quantité sous la racine vérifie :
\[
\forall x \in \R, \quad u(x) = x^2 + \e^{2x} \ge \e^{2x} > 0.
\]
La fonction $u$ est dérivable sur $\R$ (somme de fonctions dérivables) et prend des valeurs strictement positives sur $\R$.
Par composition avec la fonction racine carrée $t \mapsto \sqrt{t}$ (strictement dérivable sur $]0\,;\,+\infty[$), la fonction $f = \sqrt{u}$ est \textbf{dérivable sur $\R$}. \hfill \textbf{[0{,}5 point]}

Calculons sa dérivée pour tout réel $x$ :
\[
f'(x) = \frac{u'(x)}{2\sqrt{u(x)}} = \frac{2x + 2\e^{2x}}{2\sqrt{x^2 + \e^{2x}}} = \frac{2(x + \e^{2x})}{2\sqrt{x^2 + \e^{2x}}} = \frac{x + \e^{2x}}{\sqrt{x^2 + \e^{2x}}}.
\]
On identifie immédiatement $g(x) = x + \e^{2x}$ au numérateur, d'où :
\[
\mathbf{f'(x) = \frac{g(x)}{\sqrt{x^2 + \e^{2x}}}}.
\] \hfill \textbf{[1 point]}

\medskip
\textbf{2. (a) Coordonnées du point $A$ minimisant la distance à l'origine :}

Soit $M(x\,;\,y)$ un point quelconque de la courbe $\mathcal{C}_{\exp}$. Ses coordonnées sont $(x\,;\,\e^x)$.
La distance euclidienne du point $M$ à l'origine $O(0\,;\,0)$ dans le repère orthonormé est :
\[
OM = \sqrt{(x - 0)^2 + (\e^x - 0)^2} = \sqrt{x^2 + (\e^x)^2} = \sqrt{x^2 + \e^{2x}} = f(x).
\]
Pour trouver le point minimisant la distance $OM$, il suffit de déterminer le minimum de la fonction $f$.
Pour tout $x \in \R$, le dénominateur $\sqrt{x^2 + \e^{2x}}$ est strictement positif.
Le signe de la fonction dérivée $f'(x)$ est donc \textbf{strictement identique au signe de $g(x)$} :
\begin{itemize}[leftmargin=*]
    \item Sur l'intervalle $]-\infty\,;\,\alpha[$, $f'(x) < 0$, donc la fonction distance $f$ est strictement décroissante ;
    \item En $x = \alpha$, $f'(\alpha) = 0$ ;
    \item Sur l'intervalle $]\alpha\,;\,+\infty[$, $f'(x) > 0$, donc la fonction distance $f$ est strictement croissante.
\end{itemize}
La fonction $f$ atteint donc son \textbf{minimum global absolu et unique} en $x = \alpha$. \hfill \textbf{[1 point]}

Le point $A$ le plus proche de l'origine est donc le point d'abscisse $\alpha$.
Ses coordonnées sont :
\[
\mathbf{A\bigl(\alpha\,;\,\e^\alpha\bigr)}.
\]
\textit{Remarque fondamentale de simplification :}
Par définition de $\alpha$, $g(\alpha) = 0 \iff \alpha + \e^{2\alpha} = 0 \iff \e^{2\alpha} = -\alpha$.
Puisque $\e^{2\alpha} = (\e^\alpha)^2$ et que $\e^\alpha > 0$ (avec $\alpha < 0$), on en déduit par passage à la racine carrée :
\[
\e^\alpha = \sqrt{-\alpha}.
\]
Les coordonnées du point $A$ peuvent donc également s'écrire de manière remarquable sous la forme :
\[
\mathbf{A\bigl(\alpha\,;\,\sqrt{-\alpha}\bigr)}.
\] \hfill \textbf{[1 point]}

\medskip
\textbf{2. (b) Preuve de l'orthogonalité normale entre $(OA)$ et $T_A$ :}

\underline{\textbf{Méthode 1 : Par le produit des pentes (coefficients directeurs)}}
\begin{itemize}[leftmargin=*]
    \item \textbf{Pente de la droite $(OA)$ :}
    Puisque $O$ a pour coordonnées $(0\,;\,0)$ et $A$ a pour coordonnées $(\alpha\,;\,\e^\alpha)$ avec $\alpha \neq 0$ :
    \[
    m_{(OA)} = \frac{y_A - y_O}{x_A - x_O} = \frac{\e^\alpha - 0}{\alpha - 0} = \frac{\e^\alpha}{\alpha}.
    \]
    \item \textbf{Pente de la tangente $T_A$ :}
    La fonction exponentielle standard $\exp : x \mapsto \e^x$ a pour fonction dérivée elle-même : $(\e^x)' = \e^x$.
    Le coefficient directeur de la tangente $T_A$ à $\mathcal{C}_{\exp}$ au point d'abscisse $\alpha$ est :
    \[
    m_{T_A} = \exp'(\alpha) = \e^\alpha.
    \]
    \item \textbf{Produit des pentes dans le repère orthonormé :}
    Calculons le produit de ces deux coefficients directeurs :
    \[
    m_{(OA)} \times m_{T_A} = \frac{\e^\alpha}{\alpha} \times \e^\alpha = \frac{\e^{2\alpha}}{\alpha}.
    \]
    Or, d'après la relation caractéristique vérifiée par $\alpha$ issue de $g(\alpha) = 0$ :
    \[
    \alpha + \e^{2\alpha} = 0 \iff \e^{2\alpha} = -\alpha.
    \]
    En remplaçant $\e^{2\alpha}$ par $-\alpha$ au numérateur, on obtient immédiatement :
    \[
    m_{(OA)} \times m_{T_A} = \frac{-\alpha}{\alpha} = \mathbf{-1}.
    \]
    Dans un repère orthonormé, le produit des pentes étant égal à $-1$, les droites $(OA)$ et $T_A$ sont \textbf{rigoureusement perpendiculaires}. \hfill \textbf{[1{,}5 point]}
\end{itemize}

\underline{\textbf{Méthode 2 : Par le produit scalaire (Alternative vectorielle de référence)}}
\begin{itemize}[leftmargin=*]
    \item Le vecteur $\vec{OA}$ a pour coordonnées : $\vec{OA}\begin{pmatrix} \alpha \\ \e^\alpha \end{pmatrix}$.
    \item La tangente $T_A$ ayant pour coefficient directeur $\e^\alpha$, un vecteur directeur de $T_A$ est : $\vec{v}\begin{pmatrix} 1 \\ \e^\alpha \end{pmatrix}$.
    \item Dans le repère orthonormé, le produit scalaire euclidien vaut :
    \[
    \vec{OA} \cdot \vec{v} = \alpha \times 1 + \e^\alpha \times \e^\alpha = \alpha + \e^{2\alpha} = g(\alpha).
    \]
    Comme $g(\alpha) = 0$, on a $\vec{OA} \cdot \vec{v} = 0$.
    Les vecteurs non nuls $\vec{OA}$ et $\vec{v}$ sont orthogonaux, ce qui démontre avec une concision parfaite la perpendicularité de la droite $(OA)$ et de la tangente $T_A$.
\end{itemize}
\end{correctionbox}

\begin{center}
\begin{tikzpicture}[scale=1.4, cap=round, >=Latex]
    % Grille et axes orthonormés
    \draw[very thin, gray!30] (-2.5,-1) grid (2.5,3.5);
    \draw[->, thick] (-2.5,0) -- (2.8,0) node[right] {\small $x$};
    \draw[->, thick] (0,-1) -- (0,3.8) node[above] {\small $y$};
    \node[below left] at (0,0) {\small $O$};
    \draw[thick] (1,0.06) -- (1,-0.06) node[below] {\small $1$};
    \draw[thick] (0.06,1) -- (-0.06,1) node[left] {\small $1$};

    % Tracé de l'exponentielle
    \draw[domain=-2.4:1.25, samples=150, very thick, llgblue] 
        plot (\x, {exp(\x)}) 
        node[right] {\small $\mathcal{C}_{\exp} : y = \e^x$};

    % Définition des coordonnées de alpha et A
    % alpha approx -0.4263, e^alpha approx 0.6529
    \coordinate (O) at (0,0);
    \coordinate (A) at (-0.4263, 0.6529);

    % Segment [OA] (normale)
    \draw[very thick, ForestGreen] (O) -- (A) node[midway, left=2pt] {\footnotesize $\mathbf{d_{\min}}$};
    \filldraw[llgblue] (O) circle (2pt);
    \filldraw[ForestGreen] (A) circle (2.2pt) node[above left=2pt] {\small $A(\alpha\,;\,\e^\alpha)$};

    % Tangente T_A en A
    % Pente = e^alpha approx 0.6529
    % Droite : y - 0.6529 = 0.6529*(x - (-0.4263))
    % En x = -1.8 : y = 0.6529 - 0.6529*1.3737 approx -0.244
    % En x = 1.0  : y = 0.6529 + 0.6529*1.4263 approx 1.584
    \draw[thick, alertframe, domain=-1.8:1.1, samples=50] 
        plot (\x, {0.6529 + 0.6529*(\x + 0.4263)}) 
        node[below right] {\small $T_A$ (Tangente)};

    % Codage d'angle droit en A
    % Vecteur unitaire OA vers O : (0.4263, -0.6529) / 0.7797 = (0.5467, -0.8373)
    % Vecteur unitaire tangent : (1, 0.6529) / 1.1942 = (0.8374, 0.5467)
    \draw[thick, black] 
        ($ (A) + 0.2*(0.5467, -0.8373) $) -- 
        ($ (A) + 0.2*(0.5467, -0.8373) + 0.2*(0.8374, 0.5467) $) -- 
        ($ (A) + 0.2*(0.8374, 0.5467) $);

    % Lignes de rappel
    \draw[dotted, thick, gray] (-0.4263, 0) -- (A) -- (0, 0.6529);
    \node[below] at (-0.4263, 0) {\footnotesize $\alpha \approx -0{,}43$};
    \node[right] at (0, 0.6529) {\footnotesize $\e^\alpha \approx 0{,}65$};
\end{tikzpicture}
\end{center}

\begin{remarquebox}[Géométrie différentielle supérieure : Principe variationnel de la normale]
Le résultat démontré dans cet exercice est en réalité l'incarnation d'un \textbf{théorème fondamental de géométrie différentielle} valable pour toute sous-variété régulière de l'espace euclidien $\R^n$.

Soit $\mathcal{C}$ une courbe paramétrée régulière $t \mapsto M(t)$ de classe $\mathcal{C}^1$, et soit $O$ un point extérieur.
La distance euclidienne de $O$ au point courant $M(t)$ est minimale en un point $A = M(t_0)$ si et seulement si la fonction scalaire « carré de la distance » :
\[
\Phi(t) = \|\vec{OM}(t)\|^2 = \vec{OM}(t) \cdot \vec{OM}(t)
\]
admet un minimum local en $t_0$. La condition d'annulation de la dérivée (point critique) impose :
\[
\Phi'(t_0) = 2 \, \vec{OM}(t_0) \cdot \vec{M}'(t_0) = 0 \iff \vec{OA} \cdot \vec{T} = 0,
\]
où $\vec{T} = \vec{M}'(t_0)$ est le vecteur tangent à la courbe au point $A$.
Ainsi, \textbf{la direction de plus courte distance joignant un point fixe à une courbe est invariablement et nécessairement orthogonale à la tangente de la courbe en ce point de contact} (la droite $(OA)$ est la droite \textbf{normale} à la courbe en $A$). Ce résultat s'applique identiquement aux surfaces dans $\R^3$, aux géodésiques en relativité générale et à la théorie du contrôle optimal.
\end{remarquebox}

\vspace{0.4cm}

% ====================================================
% CORRIGÉ EXERCICE 4
% ====================================================
\subsection*{Exercice 4 : Inégalité arithmético-géométrique et convexité \hfill \normalfont\normalsize (2 points)}

\begin{correctionbox}
\underline{\textbf{Méthode 1 : Démonstration algébrique fondamentale (Moyennes Arithmétique et Géométrique)}}

Soient $x$ et $y$ deux nombres réels quelconques.
Posons les variables strictement positives :
\[
u = \e^x > 0 \quad \text{et} \quad v = \e^y > 0.
\]
Exprimons les deux quantités de l'énoncé en fonction de $u$ et $v$ :
\begin{enumerate}
    \item Par propriétés élémentaires de l'exponentielle et de la racine carrée :
    \[
    \e^{\frac{x+y}{2}} = \e^{\frac{x}{2}} \times \e^{\frac{y}{2}} = \sqrt{\e^x} \times \sqrt{\e^y} = \sqrt{\e^x \e^y} = \sqrt{uv}.
    \]
    La quantité $\e^{\frac{x+y}{2}}$ représente la \textbf{moyenne géométrique} des réels $u$ et $v$.
    \item D'autre part :
    \[
    \frac{\e^x + \e^y}{2} = \frac{u+v}{2}.
    \]
    Cette quantité représente la \textbf{moyenne arithmétique} des réels $u$ et $v$.
\end{enumerate}

Pour comparer ces deux quantités, formons leur différence :
\[
\frac{u + v}{2} - \sqrt{uv} = \frac{u - 2\sqrt{uv} + v}{2} = \frac{(\sqrt{u})^2 - 2\sqrt{u}\sqrt{v} + (\sqrt{v})^2}{2} = \frac{(\sqrt{u} - \sqrt{v})^2}{2}.
\]
Or, le carré de tout nombre réel est positif ou nul :
\[
\forall (u, v) \in (\R_+^*)^2, \quad (\sqrt{u} - \sqrt{v})^2 \ge 0 \implies \frac{u+v}{2} - \sqrt{uv} \ge 0 \iff \frac{u+v}{2} \ge \sqrt{uv}.
\]
En réinjectant $u = \e^x$ et $v = \e^y$, on en déduit l'inégalité fondamentale :
\[
\mathbf{\frac{\e^x + \e^y}{2} \ge \e^{\frac{x+y}{2}}}.
\] \hfill \textbf{[1{,}5 point]}

\medskip
\textbf{Condition d'égalité :}
L'égalité a lieu si et seulement si la différence s'annule :
\[
\frac{(\sqrt{u} - \sqrt{v})^2}{2} = 0 \iff \sqrt{u} - \sqrt{v} = 0 \iff \sqrt{u} = \sqrt{v} \iff u = v.
\]
En revenant aux variables $x$ et $y$ :
\[
\e^x = \e^y \iff \mathbf{x = y},
\]
en vertu de la stricte injectivité de la fonction exponentielle sur $\R$.
L'inégalité est donc stricte dès que $x \neq y$. \hfill \textbf{[0{,}5 point]}
\end{correctionbox}

\begin{remarquebox}[Lien avec la stricte convexité de l'exponentielle et l'inégalité de Jensen]
D'un point de vue supérieur, cette inégalité est la manifestation directe de la \textbf{stricte convexité} de la fonction exponentielle.
En effet, la fonction $\exp : t \mapsto \e^t$ est deux fois dérivable sur $\R$ et vérifie :
\[
\forall t \in \R, \quad \exp''(t) = \e^t > 0.
\]
Sa dérivée seconde étant strictement positive, la fonction exponentielle est \textbf{strictement convexe} sur $\R$.
Par définition géométrique et analytique de la convexité (ou inégalité de Jensen au milieu avec le barycentre $\lambda = \frac{1}{2}$) :
\[
\forall (x, y) \in \R^2, \quad \exp\left(\frac{x+y}{2}\right) \le \frac{\exp(x) + \exp(y)}{2}.
\]
Géométriquement, cela signifie que sur tout segment $[x\,;\,y]$, la courbe de l'exponentielle est strictement située sous la corde joignant les points $(x\,;\,\e^x)$ et $(y\,;\,\e^y)$, le milieu de la corde ayant pour ordonnée la moyenne arithmétique $\frac{\e^x+\e^y}{2}$, supérieure à l'ordonnée de la courbe au milieu $\e^{\frac{x+y}{2}}$.
\end{remarquebox}

\begin{center}
\begin{tikzpicture}[scale=1.1, cap=round, >=Latex]
    % Grille et axes
    \draw[very thin, gray!25] (-1,-0.5) grid (4,5);
    \draw[->, thick] (-1,0) -- (4.3,0) node[right] {\small $t$};
    \draw[->, thick] (0,-0.5) -- (0,5.2) node[above] {\small $y$};
    \node[below left] at (0,0) {\small $O$};

    % Courbe exponentielle
    \draw[domain=-0.8:2.1, samples=100, very thick, llgblue] 
        plot (\x, {0.6*exp(\x)}) 
        node[above right] {\small $y = \e^t$};

    % Points x et y
    % Choix : x = 0.2, y = 1.8
    \coordinate (P) at (0.2, {0.6*exp(0.2)});
    \coordinate (Q) at (1.8, {0.6*exp(1.8)});
    \coordinate (Mcorde) at (1.0, {0.5*(0.6*exp(0.2) + 0.6*exp(1.8))});
    \coordinate (Mcourbe) at (1.0, {0.6*exp(1.0)});

    % Corde joignant P et Q
    \draw[thick, alertframe] (P) -- (Q) node[midway, above left=2pt] {\footnotesize Corde sécante};
    \filldraw[alertframe] (P) circle (2pt) node[left] {\footnotesize $(x\,;\,\e^x)$};
    \filldraw[alertframe] (Q) circle (2pt) node[right] {\footnotesize $(y\,;\,\e^y)$};

    % Points au milieu
    \filldraw[ForestGreen] (Mcorde) circle (2.2pt) node[above left=1pt] {\footnotesize $\dfrac{\e^x+\e^y}{2}$};
    \filldraw[llgblue] (Mcourbe) circle (2.2pt) node[below right=1pt] {\footnotesize $\e^{\frac{x+y}{2}}$};

    % Segment vertical montrant l'écart de convexité
    \draw[dashed, thick, ForestGreen] (1,0) -- (Mcorde);
    \draw[dotted, thick, gray] (P) -- (0.2, 0) node[below] {\small $x$};
    \draw[dotted, thick, gray] (Q) -- (1.8, 0) node[below] {\small $y$};
    \node[below] at (1,0) {\small $\frac{x+y}{2}$};
\end{tikzpicture}
\end{center}

\end{document}
